\scanpage{44}
% archive sheet number: 43
% upper right corner of the sheet is torn.
$\Pi$ une $n$-carte \inserted{comb.} de moquette $R=\operatorname{Rep}(\Pi)$, $G_n$-ensemble (où
\[
G_n=\{\sigma_0,\ldots,\sigma_n\mid\sigma_0^2=\cdots=\sigma_n^2=1,\ \sigma_i\sigma_j=\sigma_j\sigma_i\text{ si }j\geq i+\unsure{2}\}
\]

Une \uline{représentation géom} $\unsure{\varphi|\phi}$ de $\Pi$ dans \deleted{\unsure{de rg $n+1$}} un espace affine $E$ (sur \unsure{anneau} $k$ disons) est la donnée d'\deleted{\unsure{une hom.}}application
\diagram{(1)\quad\begin{tikzcd}[ampersand replacement=\&]
R\arrow[r,"\varphi_R"]\&\operatorname{Drap\,max}(E)\arrow[d,equal]\\
\&\{(E_0\subset E_1\cdots\subset E_i\subset\cdots\subset E_n\subset E)\}
\end{tikzcd}}
\inserted{$E_i$: ss-espace affine de $E$ de rg $i$}

telle que pour $r\in R$, $0\leq i\leq n$, on ait
\[
\left\{\begin{aligned}
\pi_j(\varphi_R(\sigma_i r))&=\pi_j\varphi_R(r)\quad\text{si}\\
0\leq i,j&\leq n,\qquad i\ne j
\end{aligned}\right.
\tag*{(2)}
\]
où $\pi_j(d)$ désigne la composante de dim $j$ d'un drapeau $d$. Il revient au \unsure{$=$|même} \inserted{\unsure{de par.|de dire}} à dire que par passage au quotient, $\pi_j\varphi_R$ se factorise en
\diagram{(3)\quad\begin{tikzcd}[ampersand replacement=\&]
D_j(\Pi)\arrow[r,"\varphi_j"]\&\unsure{\mathrm{Gr\,aff}_j|\mathrm{Grass\,aff}_j}(E)
\end{tikzcd}\quad$(0\leq j\leq n)$}
